\section{Related work}
\label{sec:related}

\paragraph{Geometric representations and Koopman prediction.}
KAHMs learn reconstruction models in reproducing kernel Hilbert spaces
(RKHSs) \citep{kumar2024kahm,kumar2025collaborative}. We use these established
reconstruction models to construct forecasting features.
Soft state aggregation \citep{singh1994soft}, approximate state abstraction
\citep{abel2016abstraction}, and task-oriented latent transition models
\citep{gelada2019deepmdp} likewise distinguish a representation from its
predictive use. Extended dynamic mode decomposition approximates Koopman
operators on finite dictionaries \citep{williams2015edmd}; learned dictionaries
\citep{li2017dictionary} and RKHS regression bounds \citep{kostic2022regression}
provide related forecasting frameworks. Basis-invariant consistency diagnostics
\citep{haseli2023consistency}, multistep projection-error guarantees and
refinement \citep{conradie2026trustworthy}, and residual-based dictionary
learning \citep{coote2026residual} address dictionary quality. Our decisions
instead compare expected observed-output error with a declared tolerance
while preserving a common target and sampling law.

\paragraph{Closest KAHM prediction-limit analysis.}
The earlier analysis of \citet{kumar2026limits}, in its cited third version, provides
class-centered lower certificates, empirical-minimum comparisons, fitted-risk
upper bounds, and independent evaluation. It also gives an exact four-state
example and independent dynamical confirmation. We reuse its KAHM
construction and sampling principle; neither certificates for KAHM nor
independent confirmation is new here. Table~\ref{tab:prior-comparison}
states the substantive differences. Their reconstruction-score margin bounds
concern agreement with reference labels, which does not imply predictive
sufficiency.

\begin{table}[t]
\centering\small
\begin{tabular}{@{}p{0.19\linewidth}p{0.35\linewidth}p{0.40\linewidth}@{}}
\toprule
Question & \citet{kumar2026limits} & Present work \\
\midrule
Prediction target & Successor soft coordinates & Fixed observed future outputs \\
Decoder class & Spectral-norm-bounded linear coordinate predictors & General Lipschitz oracle; retained-block affine fits \\
Lower information & Class-centered variation; empirical-minimum comparator & Positive-part pair slack; normalized monotone block hierarchy \\
Representation decisions & Accuracy exclusion for a fixed coordinate map and norm limit & Simultaneous class exclusion and fitted-model acceptance; minimum feasible cost among declared alternatives \\
\bottomrule
\end{tabular}
\caption{Scope relative to arXiv:2609.32652v3. KAHM reconstruction,
independent evaluation, and standard concentration are shared foundations.
The scalar block consistency result requires the conditions in
Theorem~\ref{thm:block-hierarchy}.}
\label{tab:prior-comparison}
\end{table}

\paragraph{Pairwise bounds and Lipschitz regression.}
\citet{akash2021inverse} derive a Lipschitz-constrained mean squared error (MSE) bound
using paired residuals and target separation. Their argument is a close
antecedent to the geometric step in our witness; our contribution concerns
its conditional positive-part form, sharpness and slack, and statistically
valid stratification under the unchanged marginal risk. Finite-support
Lipschitz regression can be posed as a constrained optimization problem
\citep{gottlieb2017regression}; we use this class directly for an empirical-minimum
comparator. Model-class falsification \citep{muller2025falsification} and
minimal-empirical-risk bounds \citep{wu2026minimal} provide broader context.
The Gamma test uses local pair differences to estimate regression noise
\citep{evans2002gamma}. Lipschitz set-membership identification
\citep{milanese2004setmembership}, trajectory invalidation
\citep{jin2020invalidation}, and corruption-robust prediction ranges
\citep{blum2026reliable} concern related feasibility questions with different
sampling and error criteria.

\paragraph{Empirical minima and optimization over blocks.}
Expected sample-average optimization values and repeated-batch lower bounds
are established in stochastic programming
\citep{linderoth2006empirical,guigues2017confidence}. Optimization of empirical
$U$-statistics and the tradeoff between computation and statistical precision
are also established \citep{clemencon2016scaling}. Our block objective can be
written as the minimum of a residual-difference $U$-statistic. We do not claim
that empirical-minimum lower bounds or block averaging are new. The specific
connection is a centered hierarchy for scalar Lipschitz regression: its
size-two member is the pair witness, its expectations increase to the class
optimum under compactness, and exact examples quantify why a fixed block size
can remain inadequate. The recentering argument uses closure under adding
constants; the pair formula and scalar extension connect it to the Lipschitz
prediction problem. Calibration uses standard bounded-mean inequalities.

\paragraph{Calibration, refinement, and dependence.}
We use classical bounded-mean inequalities \citep{hoeffding1963} and empirical
Bernstein and standard-deviation bounds \citep{maurer2009empirical}. The
simultaneous testing principle follows Learn then Test
\citep{angelopoulos2025ltt}. Our group norm is established in group selection
and multiple-kernel learning \citep{bach2008group}; here its max-block dual
metric gives a comparable sensitivity budget as dictionaries grow. Soft gating
and residual fitting also have established precedents
\citep{jordan1994experts,buhlmann2003boosting}. The additional cost conclusion
comes from jointly excluding cheaper classes and accepting a fitted model.

Within-episode dependence is allowed because concentration is applied to
independent episode summaries. A chronological split of one dependent record
does not provide that independence. Forecasting bounds under other sampling
assumptions \citep{mcdonald2017forecasting} and dependence-aware conformal
analysis \citep{oliveira2024conformal} illustrate the additional conditions
needed; their guarantees are not consequences of the present theorem.
